\slajdid{07-s016}
\begin{frame}[label=07-s016]{Kako se koristi VHDL?}
\gusttekst\setlength{\parskip}{4pt}
VHDL je programski jezik opšte namene optimizovan za dizajn elektronskih kola. Kao takav, postoji mnogo tačaka u celokupnom procesu dizajna u kojima VHDL može pomoći. \par\vspace{4mm}Za specifikaciju dizajna: VHDL se može koristiti na početku, dok još dizajnirate na visokom nivou, da biste prikazali performanse i zahteve interfejsa svake komponente u velikom sistemu. Ovo je posebno korisno za velike projekte koji uključuju mnogo članova tima. Koristeći pristup dizajnu odozgo nadole, dizajner sistema može definisati interfejs za svaku komponentu u sistemu i opisati zahteve prihvatanja tih komponenti u obliku testbench-a visokog nivoa. Definicija interfejsa (obično izražena kao deklaracija VHDL entiteta) i specifikacija performansi visokog nivoa mogu se zatim preneti na druge članove tima radi dovršavanja ili usavršavanja. \par\vspace{4mm}Za unos dizajna: Unos dizajna je ona faza u kojoj se detalji sistema unose u sistem projektovanja zasnovan na računaru. U ovoj fazi možete izraziti svoj dizajn (ili delove vašeg dizajna) kao šeme (bilo na nivou ploče ili čisto funkcionalne) ili koristeći VHDL opise. Ako ćete koristiti tehnologiju sinteze, onda ćete želeti da napišete VHDL delove dizajna koristeći stil VHDL koji je prikladan za sintezu. Faza unosa dizajna može uključivati alate i metode unosa dizajna koje nisu VHDL. U mnogim slučajevima, opisi dizajna napisani u VHDL-u se kombinuju sa drugim prikazima, kao što su šeme, da bi se formirao kompletan sistem.
\end{frame}
